\documentclass{beamer}
\usetheme{Madrid}

% pdflatex
\usepackage[utf8]{inputenc}
\usepackage[T2A]{fontenc}
\usepackage[russian,english]{babel}
\usepackage{cmap}
\usepackage{babel}

\usepackage{graphicx}
\usepackage{xcolor}
\usepackage{tikz}
\usepackage{amsmath,amssymb,amsthm}
\usepackage{mathtools}
\usepackage{listings}
\usepackage{booktabs}
\usepackage{array}
\usetikzlibrary{arrows.meta,positioning,shapes.geometric,fit,backgrounds,calc,decorations.pathreplacing}

% Footline: show current section (left) and frame numbers (right)
\setbeamertemplate{footline}{%
  \leavevmode%
  \hbox{%
    \begin{beamercolorbox}[wd=.78\paperwidth,ht=2.5ex,dp=1ex,left]{author in head/foot}%
      \hspace*{1em}\usebeamerfont{footline}\insertsectionhead%
    \end{beamercolorbox}%
    \begin{beamercolorbox}[wd=.22\paperwidth,ht=2.5ex,dp=1ex,right]{date in head/foot}%
      \usebeamerfont{footline}\insertframenumber{} / \inserttotalframenumber\hspace*{1em}%
    \end{beamercolorbox}%
  }%
}

\newcommand{\parpause}[1]{\only<+->{#1\par}}
\newenvironment{definitionbox}[1][Определение]{%
  \begin{exampleblock}{#1}
}{%
  \end{exampleblock}
}

\AtBeginSection[]
{
  \begin{frame}
    \frametitle{Содержание}
    \tableofcontents[currentsection]
  \end{frame}
}

\title{Семинар 12. Хранение и обработка больших данных}
\subtitle{Распределенные вычисления}
\author{Георгий Семенов}
\institute{Институт прикладных компьютерных наук \\ Университет ИТМО}
\date{весна 2026}

\begin{document}

\frame{\titlepage}

\begin{frame}

\frametitle{Организационные моменты}

\begin{itemize}
  \item Сегодня – (14 мая, 23:59) дедлайн по домашнему заданию 
  \item Последнее пятое домашнее задание - YTsaurus и MapReduce
  \begin{itemize}
    \item Дедлайн – 28 мая, 23:59, в тот же день - последний семинар
    \item Задание без автотестов, с ручной проверкой
    \item Детальная разбалловка находится в README.MD
  \end{itemize}
  \item Что делать, чтобы закрыться?
  \begin{itemize}
    \item Делать домашние задания со штрафом 50\%, чтобы набрать 60\% от общего количества баллов на тройку (и написать мне о том, что можно проверять)
    \item Если не получается добрать баллы, то для желающих будет устный экзамен
  \end{itemize}
\end{itemize}

\end{frame}

\begin{frame}

\frametitle{Что обсуждаем сегодня?}

\begin{itemize}
  \item Что такое большие данные? (внешняя презентация)
  \item Что такое Hadoop и MapReduce-задачи? (внешняя презентация)
  \item Смотрим YTsaurus
\end{itemize}

\end{frame}

% что такое большие данные (нереляционный мир)
% хранение больших данных (DFS, файлы - какие структуры данных, Apache Parquet, Vero, ...)
% обработка больших данных (batch vs. stream, MapReduce, Spark, Flink, ...)
% Hadoop
% YTsaurus (статические и динамические таблицы)

% Протоколы аутентификации

% DIE = Distributed 

% https://www.copado.com/resources/blog/making-die-model-security-vs-the-cia-security-triad-complementary-not-competitive
% https://www.linkedin.com/pulse/cyber-security-fundamentals-cia-triad-aaa-bio-key-international


\begin{frame}
\frametitle{Итоги}

\begin{itemize}
    \item Сегодня дедлайн по домашнему заданию 4 (Raft)
    \item Выдано последнее домашнее задание 5 (YTsaurus и MapReduce)
    \item Осталось два семинара
\end{itemize}

\end{frame}

\end{document}